%!%
% Copyright (c) 2026 mingcheng <mingcheng@outlook.com>
% 
% File: ch08-prompting-advanced.tex
% Author: mingcheng <mingcheng@outlook.com>
% File Created: 2026-03-27 15:50:02
% 
% Modified By: mingcheng <mingcheng@outlook.com>
% Last Modified: 2026-03-29 07:50:31
%%

% !TeX root = ../prompt-engineering-guide-for-beginners.tex

\chapter{提示工程高阶：前沿探索与自动化}
\label{ch:prompting-advanced}

这一章介绍的技巧相对前沿，有些涉及较复杂的概念。
如果你刚开始学习提示工程，完全可以先跳过这一章，等前面的基础和进阶技巧用熟了再回来看。
把这章当作一个「技术全景图」，了解一下提示工程还有哪些更高级的玩法。

\section{思维树（Tree of Thoughts, ToT）}

在第~\ref{ch:prompting-basics}~章我们学了链式思考（CoT），让 AI 沿着一条推理路径一步步思考。但有些问题不是一条直线就能解决的，中间可能有多个分叉点，每个分叉点都可以有不同的选择。

思维树就是让 AI 在每个关键节点展开多个分支，在这个过程中采用广度优先搜索（BFS）或深度优先搜索（DFS）的算法思路评估哪条路最有希望，然后继续深入探索。如果某条路走不通（死胡同），就执行回溯（Backtracking）到上一个节点，换一条路试试。

\textit{类比：下棋时在脑子里推演好几步，每一步都想几种走法，选最优的继续。AI 也可以这样。}

\textbf{适用场景：}需要「多步试错」的规划类任务、创意类任务，或者有多个可行方案需要比较优劣的决策问题。

\textbf{使用方法：}

\begin{promptbox}
我需要为一家新开的奶茶店制定开业推广方案。

请按照以下方式思考：
\begin{enumerate}
  \item 先列出 3 种不同的推广策略方向
  \item 对每种策略，分析其优缺点和预期效果
  \item 选择最有潜力的策略，深入展开具体执行方案
  \item 如果最终方案存在明显风险，回到第 1 步考虑其他方向
\end{enumerate}
\end{promptbox}

\section{自动推理并使用工具（ART）}

现在越来越多的 AI 系统不只是「说」，还能「做」。ART 让 AI 在推理过程中自动判断：「这一步我需要查一下资料」「这个数据我应该用计算器算一下」，然后自动调用相应的工具。

\textit{类比：聪明的助理在写报告时，遇到不确定的数据会自己去查资料、用计算器验算，而不是全靠脑补。}

\textbf{实际例子：}当你让 ChatGPT（联网版）回答「今天上海的天气怎么样」时，它会自动判断需要搜索实时信息，然后调用搜索工具获取结果，再把搜索结果整合到回答中。

这种能力在很多现代 AI 工具中已经内置了。比如 ChatGPT 的代码解释器可以自动运行 Python 代码来做精确计算，Kimi 和 ChatGPT 的联网版可以自动搜索最新信息。你在使用这些工具时可能已经体验过了，只是没有意识到背后的机制就是 ART。

作为用户，你需要注意两件事：

\begin{itemize}
  \item \textbf{主动启用工具功能}：有些工具需要你手动开启联网搜索或代码执行功能。遇到需要实时信息或精确计算的任务时，记得检查一下相关功能有没有打开。
  \item \textbf{明确指示 AI 使用工具}：如果你发现 AI 在「心算」一个复杂的数学问题，可以直接告诉它「请写代码来计算，不要口算」，它通常会乖乖执行。同样，如果你需要最新数据，可以提示它「请搜索最新的信息来回答」。
\end{itemize}

简单记住一个判断标准：凡是涉及精确数字、实时信息或大规模数据处理的任务，让 AI 调用工具来完成永远比让它「脑补」靠谱。

\section{自动提示工程师（Automatic Prompt Engineer, APE）}

写提示词本身也是一项可以自动化的任务。APE 的思路是：让 AI 自己生成大量不同版本的提示词，然后在测试集上评估它们的效果，自动选出表现最好的那一版。

\textit{类比：你不自己调配咖啡豆的比例，而是让机器人自动试了 100 种配方，选出最好喝的那杯。}

这个技巧目前主要用在需要大批量处理任务的场景中（比如做数据标注、内容审核等），个人用户可能用不太到。但知道这个概念很重要，因为它揭示了提示工程领域的一个趋势：\textbf{未来，很多提示词的优化工作会由 AI 自己完成}。

其实你在日常使用中已经可以做一个简化版的 APE 了。当你对一条提示词的效果不满意时，可以直接让 AI 帮你改进：「这条提示词的效果不够好，请帮我优化它，使输出更加简洁和结构化。」让 AI 自己来优化提示词，往往能给出你意想不到的改进。

\section{ReAct 框架（Reasoning + Acting）}

ReAct 是「推理」（Reasoning）和「行动」（Acting）的结合。它让 AI 交替进行思考和行动：先想一想下一步该做什么，然后执行（比如搜索、查数据），根据执行结果再想一想下一步，反复循环直到完成任务。

\textit{类比：侦探破案时不是一口气推理到底，而是想一想，去现场找证据，再想一想，再去查资料……}

ReAct 的思维过程通常是这样的：

\begin{enumerate}
  \item \textbf{思考}：根据当前信息，我应该先做什么？
  \item \textbf{行动}：执行一个具体的操作（搜索、计算、查询等）
  \item \textbf{观察}：看看行动的结果是什么
  \item \textbf{思考}：基于新的结果，下一步该做什么？
  \item 重复上述过程，直到得出最终答案
\end{enumerate}

很多基于行业标准框架（如 LangChain、AutoGPT）打造的 AI 智能体（Agent）系统就是完全基于 ReAct 框架构建底层逻辑调度的。对于普通用户来说，了解这个框架有助于你更深刻地理解为什么有些拥有“外部工具调用能力（Function Calling）”的 AI 系统能脱离纯文本限制，通过执行 API 请求来自动完成一系列复杂的多步骤任务。

\section{Reflexion（自我反思）}

Reflexion 让 AI 在完成任务之后「回头看看」自己的回答，检查是否有遗漏、错误或不合理的地方，然后自行修正。

\textit{类比：实习生写完报告后自己先通读一遍，标出不确定的地方，修改后再交给你。}

你可以在提示词中要求 AI 进行自我反思：

\begin{promptbox}
请帮我写一封商务邮件。写完之后，请你自己检查一遍，看看是否有以下问题：
\begin{itemize}
  \item 语气是否恰当
  \item 是否有语法或拼写错误
  \item 信息是否完整
  \item 有没有可能被误解的表述
\end{itemize}
如果发现问题，请直接修正并给出最终版本。
\end{promptbox}

这个技巧简单但有效。在很多场景下，加一句「请检查你的回答是否有误」就能提升输出质量。

\section{程序辅助语言模型（PAL）}

AI 在自然语言推理方面很强，但在精确计算方面容易出错。PAL 的思路是：让 AI 遇到需要计算的部分时，不要用自然语言「心算」，而是写一段代码来执行。

\textit{类比：让实习生别心算了，直接打开 Excel 用公式算，结果靠谱得多。}

在支持代码执行的 AI 工具中（比如 ChatGPT 的代码解释器），你可以这样提示：

\begin{promptbox}
我有一组销售数据（见下表）。请用 Python 代码计算每个月的同比增长率，并找出增长最快的三个月。不要心算，请写代码运行。
\end{promptbox}

PAL 的关键在于一句话：\textbf{告诉 AI 别心算，用代码}。只要你在提示词里明确要求「写代码来计算」「用程序来处理」，AI 就会切换到程序化的工作模式，输出的结果会精确得多。

除了数学计算，PAL 在这些场景中也很好用：日期推算（计算两个日期之间间隔多少个工作日）、字符串处理（批量替换文本中的格式）、数据统计（从一堆数据中找出中位数、标准差等）。这些任务让 AI 用自然语言推理容易出错，但写成代码就是小菜一碟。

\section{多模态思维链提示（Multimodal CoT）}

传统的链式思考只处理文字，但现在的 AI 模型越来越多地支持图片、图表等视觉信息。多模态 CoT 就是把视觉信息也纳入推理过程：AI 不只是读文字，还能「看图」一步步分析。

\textit{类比：实习生不只是读文字报告，还能看着图表一步步分析趋势并给出结论。}

\textbf{示例：}

\begin{promptbox}
请看这张销售趋势图。

1. 先描述图表中显示的整体趋势
2. 找出明显的拐点或异常值
3. 分析可能的原因
4. 给出你的预测和建议
\end{promptbox}

\section{基于图的提示（Graph Prompting）}

有些问题涉及复杂的实体关系，比如人物关系、组织架构、知识网络等。基于图的提示让你用节点和连线的方式描述这些关系，帮助 AI 更好地理解和推理。

\textit{类比：画一张人物关系图，让实习生顺着关系线去回答「张三和李四是什么关系」。}

\textbf{示例：}

\begin{promptbox}
以下是一个公司的汇报关系：
\begin{itemize}
  \item 张总（CEO）管理王总（CTO）和李总（CFO）
  \item 王总（CTO）管理赵经理（研发）和孙经理（测试）
  \item 李总（CFO）管理周经理（财务）和吴经理（审计）
  \item 赵经理管理开发团队的陈工和刘工
\end{itemize}

请问：陈工的直属上级是谁？如果陈工想找 CEO 反映问题，中间需要经过几层汇报？
\end{promptbox}

\section{元提示（Meta-Prompting）}

到目前为止，我们写的提示词都是针对具体任务的。元提示的思路更高一层：用一个「总调度」提示来管理和协调多个子任务。

\textit{类比：你不直接干活，而是当总导演，安排策划组、设计组、文案组各自完成任务，最后你来审核拼装。}

\textbf{示例：}

\begin{promptbox}
你是一个任务调度器。我需要完成以下项目：为一款新上线的健身 App 制定上线推广方案。

请按照以下步骤进行：
\begin{enumerate}
  \item 先以「市场分析师」的身份，分析当前健身 App 市场的竞争格局
  \item 再以「品牌策划师」的身份，确定我们的差异化定位和核心卖点
  \item 然后以「新媒体运营」的身份，制定各渠道的具体推广计划
  \item 最后，汇总以上分析，输出一份完整的推广方案提纲
\end{enumerate}
\end{promptbox}

元提示非常适合处理需要多种专业视角协作的复杂任务。它本质上是把「角色扮演」和「提示链」结合起来，再加上一个统筹管理的层级。

\bigskip

以上就是提示工程的高阶技巧全景。这些技巧中有一些已经在主流 AI 工具中有了直接的应用，有一些还停留在学术研究阶段。但了解它们的存在，能帮你在遇到复杂问题时找到更多的思路和灵感。

\begin{tipbox}
这一章的绝大部分技巧属于「知道有这回事」就够了的层级。日常使用中，前两章的基础和进阶技巧已经能覆盖绝大多数场景。只有当你遇到普通方法解决不了的复杂问题时，再回来翻阅本章寻找灵感。
\end{tipbox}

方法论讲完了，接下来我们直接上干货。下一章会提供一批拿来即用的提示词模板，覆盖总结、翻译、写作、代码等高频场景。
